\scanpage{14}
% archive sheet number: 13
sont vérifiées. (Lorsque $X$ est également lisse, c'est contenu dans XVI 3.7 ; le cas général - sans même supposer $Y$ plat sur $S$, d'ailleurs - s'en déduit par \deleted{\unsure{un|du}} dévissage facile, et il est traité d'ailleurs dans SGA 2 XIV 1, conjonction de 1.10 et de 1.18.) La classe $\handwritten{\phi}_f$ de (4.5) peut alors s'interpréter aussi comme un élément \deleted{\unsure{de|du}}
\[
\deleted{\handwritten{\phi}_f}\handwritten{\gamma_{X/Y}\in}\mathrm{H}^{\handwritten{2}(d-d')}_{\inserted{\handwritten{X}}}(\handwritten{\unsure{Y|X}},A_{\handwritten{Y}}(d-d')).
\tag*{(4.8)}
\]
Cette classe s'appelle la \uline{classe de cohomologie à supports dans} $X$ \uline{de} $X$ \uline{dans} $Y$, et son image
\[
\handwritten{\gamma}_{X/Y}\handwritten{\in}\mathrm{H}^{\handwritten{2}(d-d')}(Y,A_Y(d-d'))
\tag*{(4.9)}
\]
\uline{la classe de cohomologie} (sous-entendue: à supports quelconques) \uline{de} $X$ \uline{dans} $Y$. Le formalisme des classes de cohomologie\inserted{\handwritten{(4.8) (4.9)}} ainsi \deleted{associées}\inserted{\handwritten{obtenues}} \deleted{aux} (une fois la définition étendue par linéarité à des cycles algébriques relatifs sur $Y$) sera étudié systématiquement dans SGA 5 IV. Notre propos ici est d\unsure{è|é|e}\deleted{n} donner \deleted{certaines pr} les propriétés les plus élémentaires de la classe locale (4.6), en attachant notre attention surtout au cas où $X$ est également lisse sur $S$.

Voici une liste de propriétés à expliciter qui me semble raisonnable\deleted{\unsure{;|:}} (où on pose $\deleted{p}\handwritten{\delta}=d-d'$)\deleted{;}:

 a) Si $X$ est lisse sur $S$, et si $G$ est localement\inserted{(ét)} sur $Y$ isomorphe à l'image inverse d'un complexe sur $S$,\inserted{\handwritten{p.ex. $G=(A)_Y$,}} alors (4.4) est un isomorphisme. (C'est essentiellement contenu dans XVI 3.7).

 b) Si $\deleted{p}\handwritten{\delta}=0$ i.e. $f$ un isomorphisme, (4.4) est l'identité i.e. $\handwritten{\phi}_f=1$ dans (4.5) ou (4.6).

\tikz[baseline=(c.base)]{\node (c) {c)};\handwritten{\draw (c.center) ellipse (.4cm and .5cm);}} Si $\deleted{p}\handwritten{\delta}=1$, de sorte que $X$ est un diviseur relatif sur $Y/S$, alors $\handwritten{\phi}_f$ de (4.5) i.e. (4.8), et a fortiori de (4.6) ou (4.9), est donné par la théorie de Kummer comme d'habitude.

 d) Compatibi\unsure{m|l}ité de $\handwritten{\phi}_f$ avec \uline{transitivité} et \uline{intersections} (\unsure{ces|ves} deux formules sont essentiellement identiques), dans le cas \uline{lisse}\handwritten{\unsure{,|;}}
